\subsection{Пул системы хранения: добавление и удаление узлов}

Пул (trusted storage pool) формируется командой \code{gluster peer probe}, которая выполняется на одном из узлов. На \code{node1} в пул добавлены \code{node2} и \code{node3}:
\begin{Verbatim}
sudo gluster peer probe node2
sudo gluster peer probe node3
sudo gluster peer status
sudo gluster pool list
\end{Verbatim}

Результат показан на рисунках~\ref{fig:probe}--\ref{fig:status2}: все узлы находятся в состоянии \code{Connected}, список пула одинаков на всех узлах.

\fig{3.1. Собрать пул.png}{Добавление узлов node2 и node3 в пул}{fig:probe}
\fig{3.2. Проверить состояние.png}{Состояние пула на node1}{fig:status1}
\fig{3.2. Проверить состояние (на node2).png}{Состояние пула на node2}{fig:status2}

Для удаления узла из пула используется команда \code{gluster peer detach}. Узел можно удалить только в том случае, если на нём нет бриков существующих томов. Узел \code{node3} удалён из пула (рисунок~\ref{fig:detach}), после чего добавлен обратно командой \code{peer probe} (рисунок~\ref{fig:reprobe}):
\begin{Verbatim}
sudo gluster peer detach node3
sudo gluster pool list
sudo gluster peer probe node3
sudo gluster pool list
\end{Verbatim}

\fig{3.3. Удалить узел из пула.png}{Удаление узла node3 из пула}{fig:detach}
\fig{3.4. Вернуть узел обратно.png}{Повторное добавление узла node3 в пул}{fig:reprobe}
